\documentclass[notes,11pt,aspectratio=169]{beamer}

\usepackage{amsmath}
\usepackage{amssymb}
\usepackage{amsfonts}
\usepackage{booktabs}
\usepackage{threeparttable}
\usepackage{graphicx}
\usepackage{adjustbox}
\usepackage{xspace}
\usepackage{appendixnumberbeamer}

\usepackage[
  backend=biber,
  style=authoryear,
  maxcitenames=2,
  mincitenames=1,
  maxbibnames=10,
  minbibnames=1,
  uniquelist=false,
  uniquename=false
]{biblatex}

\let\oldparencite\parencite
\renewcommand{\parencite}[2][]{%
  \footnotesize {\tiny \oldparencite[#1]{#2}} \normalsize %
}

\renewcommand*{\bibfont}{\footnotesize}

\addbibresource{causal_effect_schools.bib}

\usetheme[progressbar=frametitle]{metropolis}

\newcommand{\blue}[1]{\textcolor{blue}{#1}}
\newcommand{\red}[1]{\textcolor{red}{#1}}
\newcommand{\green}[1]{\textcolor{teal}{#1}}

\newenvironment{wideitemize}{%
  \begin{minipage}{0.90\textwidth}
  \begin{itemize}
  \addtolength{\itemsep}{1em}
}{%
  \end{itemize}
  \end{minipage}
}

\newenvironment{cwideitemize}{%
  \begin{center}
  \begin{minipage}{0.90\textwidth}
  \begin{itemize}
  \addtolength{\itemsep}{1em}
}{%
  \end{itemize}
  \end{minipage}
  \end{center}
}

\title{\texorpdfstring{\centering{How Much Do High Schools Matter?\\[0.35em]
Causal Effects on Achievement and Major Choice} \\ }{How Much Do High Schools Matter? Causal Effects on Achievement and Major Choice}}
\author[BEAM]{%
  \texorpdfstring{%
    \begin{columns}
      \column{.3\linewidth}
      \centering
      \column{.3\linewidth}
      \centering
      Bruno E. Aravena-Magüida\\ UChicago
      \column{.3\linewidth}
      \centering
    \end{columns}
  }
  {BEAM}
}
\date{\vspace{20pt}\hfill April 2026}

\begin{document}

\maketitle

\section{Motivation}

\begin{frame}{Question}
\begin{cwideitemize}
  \item How much do high schools matter for students' later academic achievement and higher-education major choice?

  \item The economics literature has strong evidence on school effects on test scores, but much less on whether schools affect \textit{field of study}. \parencite{angrist_stand_2016,dobbie_medium-term_2015,angrist_leveraging_2017,angrist_credible_2024}

  \item One empirical challenge is that ``school attended'' is a high-dimensional treatment: one treatment for each high school.
\end{cwideitemize}
\end{frame}

\begin{frame}{Original Problem: Many Treatments}
\begin{wideitemize}
  \item My original goal was to estimate idiosyncratic effects of each school.

  \item A fully flexible treatment would include one indicator for each high school:
  \[
    D_i^{\mathrm{full}}
    =
    \left(
    \mathbf{1}\{S_i=1\},
    \mathbf{1}\{S_i=2\},
    \ldots,
    \mathbf{1}\{S_i=J\}
    \right).
  \]

  \item The lottery shifts students across schools, but generally does not identify a separate causal effect for every school.
\end{wideitemize}
\end{frame}

\begin{frame}{Original Problem: Common Fallback}
\begin{wideitemize}
  \item Suppose all students rank three schools \(A \succ B \succ C\).

  \item Suppose they all have the same guaranteed fallback school \(G\).

  \item Then lottery-induced comparisons can be interpreted relative to the same counterfactual:
  \[
    \beta_{AG}, \qquad \beta_{BG}, \qquad \beta_{CG}.
  \]

  \item Because the comparison school is common, these objects are comparable across target schools.
\end{wideitemize}
\end{frame}

\begin{frame}{Original Problem: Different Fallbacks}
\begin{wideitemize}
  \item In practice, students have different fallback schools:
  \[
    G_i \in \{G^1, G^2, \ldots, G^K\}.
  \]

  \item The relevant objects become target-school-by-fallback effects:
  \[
    \beta_{AG^1}, \beta_{AG^2}, \ldots, \beta_{CG^K}.
  \]

  \item A gain from \(G^1\) to \(A\) is not automatically comparable to a gain from \(G^7\) to \(B\).

  \item This motivates projecting schools into a scalar value index instead of trying to recover every pairwise school contrast.
\end{wideitemize}
\end{frame}

\begin{frame}{Collapsing to meaningful margins: \textcite{kirkeboen_field_2016}}
\begin{wideitemize}
  \item \textcite{kirkeboen_field_2016} study field-of-study effects using centralized assignment in Norway. 

  \item Because applicants have different next-best alternatives, comparing each program would put us in a similar situation than the one described before.
  
 \item Their solution: collapse programs into broad fields and estimate next-best effects across fields. 
 
 \item The cost of this approach is that it requires ignoring effects(variation) coming from differences in institution in their main specification.

\end{wideitemize}
\end{frame}

\begin{frame}{Collapsing to meaningful margins: Next-Best Effects}
\begin{figure}
  \centering
  \includegraphics[width=1.2\linewidth,height=1.2\textheight,keepaspectratio,angle=270]{C:/Users/xd-br/Desktop/PhD/Research/causal_schools/writing/spring_update/magne_table.png}
\end{figure}
\vspace{-0.8em}
\begin{center}
  \tiny Source: \oldparencite{kirkeboen_field_2016}.
\end{center}
\end{frame}

\begin{frame}{Main Pivot}
\begin{wideitemize}
  \item Instead of estimating one causal effect for each school, collapse school attendance into a scalar school-value treatment.

  \item For each school \(s\), construct an observational school value \(\beta_s\).

  \item A student's treatment is the value of the school they attend:
  \[
    D_i = \beta_{S_i}.
  \]

  \item The IV asks whether lottery-induced movement toward higher-value schools (in a given metric) causally improves outcomes.
  \item Since the variation comes from lottery offers, the effects tells us about the causal effects of schools. But at the same time, these results are mediated by a certain projection, which might nor not be relevant. 
\end{wideitemize}
\end{frame}

%\begin{frame}{Interpretation}
%\begin{wideitemize}
%  \item The object is not the effect of attending a particular school.

%  \item The parameter summarizes the causal return to moving along a pre-specified school-value index.

%  \item In words:

%  \vspace{0.4em}
%  \begin{center}
%    \textit{When lottery offers induce students to attend schools with higher observational value, do their outcomes improve?}
%  \end{center}
%\end{wideitemize}
%\end{frame}

\section{Data and School Values}

\begin{frame}{Data}
\begin{wideitemize}
  \item Universe: students observed in regular grade 8 in Chile between 2017 and 2020.

  \item Approx. 945,000 unique students.

  \item Linked records:
  \begin{itemize}
    \item grade-4 SIMCE math and language scores;
    \item grade-8 cohort, gender, age, and home comuna;
    \item school enrollment histories after grade 8;
    \item PSU/PAES math and language scores;
    \item higher-education enrollment and STEM field indicators.
  \end{itemize}
\end{wideitemize}
\end{frame}

\begin{frame}{Assigning a High School}
\begin{wideitemize}
  \item For each student, I track school enrollment during the four years after grade 8.

  \item Main school assignment:
  \[
    S_i = \text{school code attended for the most post-grade-8 years}.
  \]

  \item Ties are broken toward the most recent school.

  \item This maps each student in the broad grade-8 universe to one high-school code when available.
\end{wideitemize}
\end{frame}

\begin{frame}{Raw and Adjusted School Values}
\small
\textbf{Raw school mean}
\[
\beta_s^{\mathrm{raw}}
=
\frac{1}{N_s}
\sum_{i:S_i=s}Y_i.
\]

\vspace{0.5em}
\textbf{Adjusted observational value-added}
\[
Y_i
=
\mu
+
\beta_{S_i}
+
\sum_{p=1}^{3}\lambda_{Mp}M_i^p
+
\sum_{p=1}^{3}\lambda_{Lp}L_i^p
+
\gamma_{\mathrm{cohort}(i)}
+
\delta_{\mathrm{gender}(i)}
+
\rho_{\mathrm{age}(i)}
+
\kappa_{\mathrm{comuna}(i)}
+
\varepsilon_i .
\]

\vspace{0.3em}
\begin{itemize}
  \item \footnotesize \(\beta_s\) is the school fixed effect. I use this as an observational, selection-contaminated measure of school value. \parencite{chetty_measuring_2014,angrist_leveraging_2017}
  \item \footnotesize For these results to be causal, we would need to assume that the controls remove all selection on unobservables that is correlated with the school attended. \textit{(CIA)}
  \item \footnotesize I have also a noisy measure of family income (G4), attendance and performance within schools (G 1-8). Some signal there but skeptical. 
\end{itemize}


\end{frame}

\begin{frame}{Adjustment Matters and Patterns Are Sensible}
\begin{figure}
  \centering
  \includegraphics[width=0.82\linewidth]{C:/Users/xd-br/Desktop/PhD/Research/causal_schools/output/figures/school_rbd_observational_values/math_values_by_region_raw_adjusted.png}
\end{figure}
\end{frame}

\section{IV Design}

\begin{frame}{Lottery Sample}
\begin{wideitemize}
  \item I use students who enter the grade-9 centralized assignment mechanism \textit{(SAE)} on time (i.e.  same year they are observed in grade 8). 

  \item I restrict to students with non-trivial assignment risk:
  \[
    \max_{s\in\mathcal{S}} p_{is}<1,
  \]
  where \(p_{is}\) is the simulated prob. that student \(i\) is assigned to school \(s\).

  \item This implies going from $\approx$ 320,000 students that enter SAE, to $\approx$ 165,000 students with assignment risk. This is roughly 1/6 of students in the universe having assignment risk. 

  \item The lottery offer provides exogenous variation in the value of the school attended, conditional on assignment-probability controls.
\end{wideitemize}
\end{frame}

\begin{frame}{Supported School Set}
\begin{wideitemize}
  \item Full probability controls are high-dimensional. Concretely 2 per school: one for the probability and one for the zero-probability indicator.

  \item I define the supported lottery school set as schools with positive assignment probability for at least 100 at-risk students:
  \[
  \mathcal{S}^{100}
  =
  \left\{
  s:
  \sum_i \mathbf{1}\{p_{is}>0\}
  \geq 100
  \right\}.
  \]

  \item This keeps the current design computationally feasible while covering a large share of students.
\end{wideitemize}
\end{frame}

\begin{frame}{Support Tradeoff}
\centering
\tiny
\begin{tabular}{llrrr}
\toprule
Threshold students at risk & School group & Schools & Students & Share \\
\midrule
100 & SAE \(<100\) & 1,491 & 315,261 & 33.4\% \\
100 & SAE \(\geq100\) & 1,011 & 493,955 & 52.3\% \\
100 & Particular Pagado & 497 & 87,373 & 9.3\% \\
100 & Other non-SAE schools & 1,134 & 47,771 & 5.1\% \\
\midrule
50 & SAE \(<50\) & 940 & 174,922 & 18.5\% \\
50 & SAE \(\geq50\) & 1,562 & 634,294 & 67.2\% \\
50 & Particular Pagado & 497 & 87,373 & 9.3\% \\
50 & Other non-SAE schools & 1,134 & 47,771 & 5.1\% \\
\midrule
25 & SAE \(<25\) & 530 & 89,266 & 9.5\% \\
25 & SAE \(\geq25\) & 1,972 & 719,950 & 76.2\% \\
25 & Particular Pagado & 497 & 87,373 & 9.3\% \\
25 & Other non-SAE schools & 1,134 & 47,771 & 5.1\% \\
\bottomrule
\end{tabular}
\end{frame}

\begin{frame}{Using the Value-Added Metric: Treatment and Instrument}
\small
The endogenous school-value treatment is
\[
D_i
=
\beta_{S_i}.
\]

The excluded instrument is
\[
Z_i
=
\mathbf{1}\{O_i\in\mathcal{S}^{100}\}
\beta_{O_i}.
\]

\vspace{0.5em}
\begin{wideitemize}
  \item \(S_i\): school attended after grade 8.
  \item \(O_i\): first-round lottery offer.
  \item Offers outside \(\mathcal{S}^{100}\) are assigned zero in the excluded instrument.
\end{wideitemize}
\end{frame}

\begin{frame}{2SLS System}
\small
\[
D_i
=
\pi_0
+
\pi_1 Z_i
+
P_i'\psi_D
+
X_i'\lambda_D
+
u_i .
\]

\[
Y_i
=
\alpha
+
\theta \widehat{D}_i
+
P_i'\psi_Y
+
X_i'\lambda_Y
+
\varepsilon_i .
\]

\vspace{0.6em}
\begin{wideitemize}
  \item \(P_i\): assignment-probability controls for schools in \(\mathcal{S}^{100}\), including probabilities and zero-probability indicators.
  \item \(X_i\): grade-8 cohort, gender, grade-8 age, fourth-grade math score, and fourth-grade language score. Probably moving to higher-order polynomials.
  \item \(\theta\): causal return to a lottery-induced increase in the scalar school-value index.
\end{wideitemize}
\end{frame}

\section{Preliminary Results}


\begin{frame}{How to Read \(\theta\)}
\begin{wideitemize}
  \item The general interpretation is that a student induced by a lottery offer to attend a school (a complier) with a 1-unit higher value in the respective metric is expected to have a $\theta$-unit higher outcome.

  \item Here, \(\theta\) compares causal gains to the observational school-value metric.
   If \(\theta=1\), then a one-unit difference in observational school value maps one-for-one into a causal effect. \ This idea is formalized and a test for it is proposed in \textcite{angrist_leveraging_2017}.

\end{wideitemize}
\end{frame}

\begin{frame}{Main Results}
\centering
\scriptsize
\begin{tabular}{lcccccc}
\toprule
 & \multicolumn{2}{c}{Math VA} & \multicolumn{2}{c}{Language VA} & \multicolumn{2}{c}{STEM VA} \\
\cmidrule(lr){2-3} \cmidrule(lr){4-5} \cmidrule(lr){6-7}
 & Unadj. & Adj. & Unadj. & Adj. & Unadj. & Adj. \\
\midrule
\(\theta\) & 0.208 & 0.431 & 0.167 & 0.393 & 0.345 & 0.584 \\
SE & 0.033 & 0.057 & 0.032 & 0.062 & 0.046 & 0.068 \\
\midrule
First-stage coef. & 0.455 & 0.457 & 0.439 & 0.435 & 0.471 & 0.455 \\
First-stage F & 4399 & 3656 & 5481 & 4899 & 6509 & 5086 \\
N & 105,902 & 105,893 & 105,902 & 105,894 & 135,588 & 135,588 \\
\bottomrule
\end{tabular}

\vspace{0.7em}
\footnotesize
All \(p\)-values are below 0.001.
\end{frame}



\begin{frame}{Interpretation}
\begin{wideitemize}

%  \item Current estimates are positive and substantial, but below one.

%  \item Adjusted metrics perform better than raw means, consistent with the idea that controls remove part of the selection component.

\item VA metrics are predictive of lottery-induced causal effects across specifications. 
\item Results are metric-specific, but also show that exogenous changes in school assignments do change outcomes for individuals. 

\item To my knowledge, results are among the first to show that school assignment lotteries induce changes in higher education major choice, which is an important contribution in itself.

\end{wideitemize}
\end{frame}

\begin{frame}{STEM Enrollment by Gender}
\centering
\scriptsize
\begin{tabular}{lcccccc}
\toprule
 & \multicolumn{3}{c}{STEM VA} & \multicolumn{3}{c}{STEM gender gap reduc. VA} \\
\cmidrule(lr){2-4} \cmidrule(lr){5-7}
 & All & Boys & Girls & All & Boys & Girls \\
\midrule
\(\theta\) & 0.584 & 0.535 & 0.695 & -0.254 & -0.477 & 0.000 \\
SE & 0.068 & 0.088 & 0.097 & 0.078 & 0.122 & 0.087 \\
p-value & \(<0.001\) & \(<0.001\) & \(<0.001\) & 0.001 & \(<0.001\) & 0.998 \\
\midrule
First-stage coef. & 0.455 & 0.460 & 0.443 & 0.443 & 0.457 & 0.429 \\
First-stage F & 5086 & 3198 & 2003 & 4954 & 2810 & 2154 \\
N & 135,588 & 67,221 & 68,367 & 125,672 & 63,669 & 62,003 \\
\bottomrule
\end{tabular}
\end{frame}

\begin{frame}{STEM Gender Interpretation}
\begin{wideitemize}
  \item Schools with higher adjusted STEM VA raise STEM enrollment for both boys and girls.

  \item The point estimate is larger for girls than boys, but the difference is not statistically significant in the current separate-sample comparison.

  \item The gender-gap-reduction VA is defined as female minus male STEM value-added.

  \item Preliminary evidence suggests that schools with higher gender-gap-reduction VA reduce the male-favored gap mainly through lower boys' STEM enrollment, rather than clear increases for girls.
\end{wideitemize}
\end{frame}


\begin{frame}{School Funding Per Person}
\begin{wideitemize}
  \item I also apply the same scalar-treatment framework to school resources.

  \item For each school, I construct public funding per student from MINEDUC subsidy payments to schools.

  \item I use two school-level measures:
  \begin{itemize}
    \item average public funding per student over 2017--2021;
    \item growth in public funding per student from 2017--2018 to 2019--2021.
  \end{itemize}

  \item Both measures are expressed in millions of 2021 pesos per student, roughly thousands of 2021 USD per student.
\end{wideitemize}
\end{frame}

\begin{frame}{Funding Per Person Results}
\centering
\scriptsize
\begin{tabular}{lcccccc}
\toprule
 & \multicolumn{3}{c}{Funding per student} & \multicolumn{3}{c}{Growth in funding per student} \\
\cmidrule(lr){2-4} \cmidrule(lr){5-7}
 & Math & Language & STEM & Math & Language & STEM \\
\midrule
\(\theta\) & -0.022 & -0.041 & 0.029 & -0.020 & -0.001 & 0.031 \\
SE & 0.027 & 0.028 & 0.014 & 0.048 & 0.050 & 0.021 \\
p-value & 0.416 & 0.144 & 0.042 & 0.685 & 0.979 & 0.132 \\
\midrule
First-stage coef. & 0.339 & 0.339 & 0.268 & 0.398 & 0.398 & 0.343 \\
First-stage F & 73.6 & 73.6 & 53.4 & 139.7 & 139.7 & 127.7 \\
N & 97,303 & 97,303 & 123,953 & 96,501 & 96,501 & 123,052 \\
\bottomrule
\end{tabular}

\vspace{0.7em}
\footnotesize
Treatments are measured in millions of 2021 pesos per student.
\end{frame}

\begin{frame}{Funding Interpretation}
\begin{wideitemize}
  \item The expenditure-based scalar metric is not consistently predictive of lottery-induced gains in test scores or STEM enrollment.

  \item The first stages are strong, so the null-ish reduced relationship is not coming from a weak instrument.
  However, the first-stage is much weaker than in the previous specifications.

  \item \textbf{What these metrics capture matters for the results:} public funding for these schools is largely determined by enrollment and thus should be constant per student within schools.
  \item The non-enrolmment component of funding tends to target specific student populations, so higher funding per student is not automatically a pure resource shock.

\end{wideitemize}
\end{frame}

\begin{frame}{Takeaways}
\small
\begin{wideitemize}
  \item The scalar school-value approach makes the lottery variation tractable without estimating one effect per school.

  \item Observational school VA predicts lottery-induced causal gains in achievement and STEM enrollment.

  \item Adjusted VA performs better than raw means.

  \item The current design suggests schools affect higher-education major choice, not only test scores.

%  \item Funding per student has strong first stages, but does not currently predict causal gains as clearly as the outcome-based VA metrics.
\end{wideitemize}
\end{frame}

\begin{frame}{Next Short-Term Steps}
\small
\begin{wideitemize}
%  \item Formalize the estimand and support restriction in the write-up.
  \item Further clean, impute and include family income in the VA metrics. 

  \item Explore run-time of the \(50\)- and \(25\)-student support thresholds.

  \item Test gender differences using pooled interaction specifications.

%  \item Refine the interpretation of the funding exercise and separate formula-driven funding from components closer to discretionary resources.
\end{wideitemize}
\end{frame}

\begin{frame}{Extension: Do students moved into STEM programs actually do better?}
  \begin{wideitemize}
    
    \item Discretize the STEM VA metric and characterize compliers using \textcite{abadie_semiparametric_2003}.
    
    \item If I had data, I can run a regression of wages on STEM enrollment --and controls-- for the (school-lottery) compliers above i.e. 
    using the $\kappa$ weights from the \textcite{abadie_semiparametric_2003} characterization.   
    
    \item This might be possible in 5 years or so since these individuals are just starting college.  An alternative is to look at intermediate outcomes like college grades or persistence in STEM majors.
    
    \item Another alternative is to look at similar populations in STEM programs in the past, under a version of Conditional Effects Ignorability (CEI) \parencite{angrist_extrapolate-ing_2010}.  


  \end{wideitemize}


\end{frame}


\begin{frame}{Distribution of STEM School VA}
\begin{figure}
  \centering
  \includegraphics[width=0.82\linewidth,height=0.72\textheight,keepaspectratio]{C:/Users/xd-br/Desktop/PhD/Research/causal_schools/output/figures/scalar_school_value_iv/stem_adj_va_distribution_sae_gt100_private_paid_n_legend.png}
\end{figure}
\end{frame}



\begin{frame}[standout]
Thanks!
\end{frame}

\end{document}
